\documentclass[11pt]{amsart}
% packages
\usepackage[margin=1.1in]{geometry}
\usepackage{amssymb,amsthm}
\usepackage{xcolor}
\usepackage[colorlinks=true,linkcolor=blue!50!black,citecolor=blue!50!black,urlcolor=blue!50!black]{hyperref}

% theorem styles
\theoremstyle{plain}
	\newtheorem{theorem}{Theorem}
	\newtheorem{lemma}[theorem]{Lemma}
	\newtheorem{corollary}[theorem]{Corollary}
\theoremstyle{remark}
	\newtheorem{remark}[theorem]{Remark}

% shortcuts
\newcommand{\Mlt}{\mathrm{Mlt}}
\newcommand{\Stab}{\mathrm{Stab}}

\newcommand{\E}{\mathbb{E}}
\renewcommand{\Pr}{\mathbb{P}}
\newcommand{\LL}{\mathcal{L}}
\DeclareMathOperator{\sgn}{sgn}

\title{Almost all loops have multiplication group $S_n$}
\author{}
\date{October 2026 --- draft}

\begin{document}
\maketitle

\begin{abstract}
Cameron proved that for almost all quasigroups of order $n$ the left
translations generate the symmetric or the alternating group, and
H\"aggkvist and Janssen excluded the alternating group; Cameron asked
whether the same holds for loops, where the argument breaks because the
translation by the identity is trivial.  We show that it does, by counting: the
probability that the rows of a uniformly random reduced Latin square of
order $n$ generate a group not containing $A_n$ is at most
$e^{-(\log 2-o(1))n^2}$, and the group is $S_n$ with probability
$1-O(c^n)$ for some $c<1$.  The inputs are the van der Waerden and
Br\'egman bounds on the number of Latin squares, the elementary bounds
of Babai and Pyber on the orders of primitive groups, and the
H\"aggkvist--Janssen bound on Latin squares all of whose rows are even
permutations; the classification of finite simple groups is not used.  This settles the question,
raised by Cameron and recorded by Cavenagh, Greenhill and Wanless as a
consequence of their conjecture on the second row, and completes an
argument sketched by Eberhard in 2016.
\end{abstract}

\section{Introduction}

Let $Q$ be a quasigroup of order $n$ on the set $[n]=\{1,\dots,n\}$,
with Cayley table a Latin square $L$, and for $x\in Q$ let
$L_x\colon y\mapsto xy$ be the left translation, the $x$-th row of $L$
read as a permutation of $[n]$.  The \emph{left multiplication group}
$\Mlt_\lambda(Q)=\langle L_x:x\in Q\rangle$ is a transitive subgroup
of $S_n$, and the full multiplication group $\Mlt(Q)$, generated by all
left and right translations, contains it.  Cameron \cite{Cam92} showed
that for almost all quasigroups (that is, for all but a vanishing
proportion of the Latin squares of order $n$) the left translations
generate $S_n$ or $A_n$; H\"aggkvist and Janssen \cite{HJ96} then showed
that the proportion of Latin squares all of whose rows are even
permutations is at most $c^n$ for some $c<1$, so that $A_n$ occurs for
a proportion at most $c^n$ of quasigroups.  In particular $\Mlt(Q)$ is $2$-transitive, which in Smith's
character theory of quasigroups \cite{Smith} means that almost all
quasigroups have trivial character theory (``rank $2$'').

Cameron's argument is short.  A uniformly random Latin square has a
uniformly random first row, and by a theorem of {\L}uczak and Pyber
\cite{LP93} a uniformly random permutation lies in no transitive proper
subgroup of $S_n$ other than $A_n$ with probability $1-o(1)$; since the
rows generate a transitive group, that group contains $A_n$.

For \emph{loops} --- quasigroups with an identity element, whose Cayley
tables, suitably labelled, are the \emph{reduced} Latin squares, with
first row and first column in natural order --- the argument is not
available: the first row is the identity permutation, and no other row
is uniformly distributed.  Cameron asked whether almost all loops nevertheless have
multiplication group $S_n$ (see \cite{Cam15}, and \cite[\S6]{CGW08},
where Cavenagh, Greenhill and Wanless record, as Lemma~6.2, that this
follows from their conjecture on the cycle structure of the second row
of a random Latin square).  The statement is folklore in loop theory,
where it is mentioned as an unproved belief.  The purpose of this note
is to settle it, by an argument that uses no information about the
second row at all.

\begin{theorem}\label{thm:main}
Let $L$ be a uniformly random reduced Latin square of order $n$ and
$H(L)\le S_n$ the group generated by its rows.  Then
\[
\Pr\bigl[H(L)\not\supseteq A_n\bigr]\le e^{-(\log2-o(1))n^2},
\qquad\text{and}\qquad
\Pr\bigl[H(L)=S_n\bigr]\ge1-3c^n
\]
for some absolute $c<1$ and all large $n$.  Consequently almost all loops
have $\Mlt_\lambda(Q)=\Mlt(Q)=S_n$ and trivial character theory.
\end{theorem}

\noindent(The first row of a reduced square is the identity, so $H(L)$
is the left multiplication group of the loop.  For a square with
arbitrary first row $r_1$ the group that matters is $\langle r_1^{-1}r_i\rangle$,
the row group of the loop whose Cayley table is obtained by relabelling
the symbols and then reordering the rows; this is set out in
\S\ref{sec:prelim}, and it is in the larger set of squares
with first row the identity that the counting is done.)

The idea, which Eberhard sketched in comments on Cameron's post
\cite{Cam15}, is that the event is about all $n-1$ non-identity rows
simultaneously, and such events are cheap to bound by counting: for a
subgroup $K\le S_n$, at most $|K|^{n-1}$ of the $L(n)/n!=(n!)^{n-1}e^{-n^2+o(n^2)}$
Latin squares with first row the identity have all their rows in $K$,
a proportion of at most $e^{n^2+o(n^2)}/[S_n:K]^{n-1}$, which is
$e^{-\delta n^2+o(n^2)}$ as soon as $[S_n:K]\ge e^{(1+\delta)n}$.  A
transitive proper subgroup lies in a maximal one, and the maximal
transitive subgroups of $S_n$ are the wreath products $S_a\wr S_b$,
$ab=n$, of which there are fewer than $n\cdot n!$, and the primitive
maximal subgroups.  The crude count disposes of $S_a\wr S_b$ for
$b\ge3$ (the index $n!/((a!)^bb!)$ is a multinomial coefficient over
$b!$, at least $3^n/n^{O(1)}$) but not of $S_{n/2}\wr S_2$
(index $\binom n{n/2}/2<e^n$) nor of $A_n$; for the primitive
subgroups one needs a bound on their orders strong enough that their
number does not matter.  We
handle every imprimitive group by a single block decomposition
(Lemma~\ref{lem:block}), which is sharper than the crude count ---
for $b=3$ the crude count gives $e^{-(\log3-1)n^2+o(n^2)}$ and for
$b=2$ nothing, while the lemma gives $e^{-(\log2-o(1))n^2}$ for every
$b$ --- and in particular disposes of $S_{n/2}\wr S_2$; the primitive groups by
the elementary order bounds of Babai and Pyber, after which their number
is controlled by counting generators; and $A_n$ by the theorem of
H\"aggkvist and Janssen.  No step uses the classification of finite
simple groups.

\section{Preliminaries}\label{sec:prelim}

\paragraph{Latin squares.}  Let $\LL_n$ be the set of Latin squares of
order $n$ on the symbol set $[n]$ and $L(n)=|\LL_n|$.  We use two
classical bounds \cite[Ch.~17]{vLW}: the lower bound
$L(n)\ge(n!)^{2n}/n^{n^2}$, a consequence of the van der Waerden
permanent inequality applied row by row, and the upper bound
$L(n)\le\prod_{k=1}^n(k!)^{n/k}$, a consequence of Br\'egman's.  In the
explicit forms we need,
\begin{equation}\label{eq:L}
\log L(n)\ge n^2\log n-2n^2,\qquad
\log L(a)\le a^2\log a-2a^2+a\bigl(\tfrac12\log^2a+2\log a+3\bigr)
\end{equation}
for all $n,a\ge1$.  The first follows from $n!\ge(n/e)^n$.  For the
second, $\log k!\le k\log k-k+\log k+1$ for $k\ge1$ gives
$(\log k!)/k\le\log k-1+(\log k+1)/k$, and summing over $k\le a$ with
$\log a!\le a\log a-a+\log a+1$, $\sum_{k\le a}(\log k)/k\le\frac12\log^2a+1$
and $\sum_{k\le a}1/k\le\log a+1$,
\[
\sum_{k=1}^a\frac{\log k!}k\le a\log a-2a+\tfrac12\log^2a+2\log a+3,
\]
which multiplied by $a$ is the claim.

\paragraph{Loops and reduced squares.}  Three sets of squares are
counted below: $\LL_n$; the squares with first row the identity,
$\LL_n^0:=\{L\in\LL_n:r_1=\mathrm{id}\}$; and the reduced squares, whose
first row and first column are the identity, $R(n)$ in number.
Relabelling the symbols of $L\in\LL_n$ by a permutation is a bijection
of $\LL_n$ (there are $n!$ of them), permuting the rows is another
($n!$ of them), and
\begin{equation}\label{eq:R}
|\LL_n^0|=\frac{L(n)}{n!},\qquad R(n)=\frac{|\LL_n^0|}{(n-1)!}=\frac{L(n)}{n!\,(n-1)!}:
\end{equation}
exactly one relabelling of the symbols (by $r_1^{-1}$, where
$r_1,\dots,r_n\in S_n$ are the rows of $L$, $r_i(c)=L(i,c)$) brings the
first row to the identity, giving the square $L^0$ with rows
$\sigma_i=r_1^{-1}r_i$; and exactly one permutation of the rows of
$L^0$ fixing the first row (which is the unique row with entry $1$ in
column $1$) brings the first column to $1,2,\dots,n$, giving a reduced
square with the same \emph{set} of rows as $L^0$.  A uniform reduced
square is the Cayley table of a uniform random loop on $[n]$ with
identity $1$, whose left translations are its rows.  Writing
\[
H(L):=\langle\sigma_2,\dots,\sigma_n\rangle=\langle r_1^{-1}r_i:2\le i\le n\rangle\le S_n,
\]
relabelling the symbols of $L$ by $\tau$ replaces each $r_i$ by $\tau r_i$
and leaves every $r_1^{-1}r_i=(\tau r_1)^{-1}(\tau r_i)$, hence $H(L)$,
unchanged; so for every subgroup $K\le S_n$ the proportion of $L\in\LL_n$
with $H(L)=K$ equals the proportion of $L^0\in\LL_n^0$ with $H(L^0)=K$,
equals the proportion of reduced squares with row group $K$, equals the
proportion of loops $Q$ of order $n$ with $\Mlt_\lambda(Q)=K$.  It
therefore suffices to prove the bounds of Theorem~\ref{thm:main} for
$H(L)$ as proportions of $\LL_n^0$, which is what we do.  Two further
facts.  First, $H(L)$ is transitive: the first column of $L$ is a
permutation, so the values $\sigma_i(1)=r_1^{-1}(L(i,1))$, $1\le i\le n$,
run through all of $[n]$.  Second, for a subgroup $K\le S_n$ the squares
$L\in\LL_n$ all of whose rows lie in $K$ are exactly the relabellings by
$\tau\in K$ of the squares $L^0\in\LL_n^0$ all of whose rows lie in $K$
(if $r_1\in K$ then $r_i\in K$ iff $\sigma_i\in K$), $|K|$ of them for
each $L^0$; dividing by $L(n)=n!\,|\LL_n^0|$,
\begin{equation}\label{eq:K}
\Pr_{\LL_n^0}\bigl[H\le K\bigr]=\frac{n!}{|K|}\,\Pr_{\LL_n}\bigl[\text{all rows in }K\bigr],
\end{equation}
where $\Pr_X$ denotes proportion in $X$.  The two sides differ by the
index of $K$, which is where a count of squares with all rows in $K$
(as in Lemma~\ref{lem:block}) and a statement about $H$ part company.

\paragraph{Group theory.}  We use three facts about subgroups of $S_n$.
A transitive subgroup not containing $A_n$ is either imprimitive, i.e.\
contained in the stabiliser $\Stab(\Pi)\cong S_a\wr S_b$ of a partition
$\Pi$ of $[n]$ into $b$ blocks of size $a$ with $1<a<n$, or primitive.
A primitive subgroup $K\not\supseteq A_n$ has order at most
\[
B:=e^{4\sqrt n\log^2n}\qquad\text{for }n\ge n_0.
\]
Indeed, if $K$ is uniprimitive (primitive but not doubly transitive),
$|K|<B$ is Babai's theorem \cite{Bab81}; if $K$ is doubly transitive,
$|K|\le\exp(\exp(c\sqrt{\log n}))$ by Babai \cite{Bab82} and
$|K|\le n^{c\log^2n}$ by Pyber \cite{Pyb93b}, and either bound is below
$B$ once $n$ is large.  The proofs of all three bounds are elementary;
in particular they do not use the classification of finite simple
groups.  Finally a group of order $m$ is generated by at most $\log_2m$
elements (adjoining to a proper subgroup an element outside it at least
doubles the order).

\section{Proof of Theorem~\ref{thm:main}}

\begin{lemma}[squares with a block structure]\label{lem:block}
Let $\Pi=\{C_1,\dots,C_b\}$ be a partition of the columns $[n]$ and
$\Pi'=\{S_1,\dots,S_b\}$ a partition of the symbols $[n]$, both into $b$
blocks of size $a$, $ab=n$.  The number of Latin squares $L\in\LL_n$
every row of which maps each block of $\Pi$ onto a block of $\Pi'$ is
exactly $T(a,b)\,L(a)^{b^2}$, where $T(a,b)\le(b!)^n$ is the number of
$n\times b$ arrays with entries in $[b]$ whose rows are permutations of
$[b]$ and in each of whose columns every symbol occurs exactly $a$
times.  (Only $\Pi'=\Pi$ is used below; the proof is the same for two
partitions.)
\end{lemma}

\begin{proof}
Each row $r_i$ of such a square induces a bijection $\pi_i\colon[b]\to[b]$
with $r_i(C_j)=S_{\pi_i(j)}$.  Fix $j,j'$ and let
$I_{jj'}=\{i:\pi_i(j)=j'\}$.  In a column $c\in C_j$ each of the $a$
symbols of $S_{j'}$ occurs exactly once, and only in rows of $I_{jj'}$,
each of which contributes exactly one symbol of $S_{j'}$ to column $c$;
hence $|I_{jj'}|=a$, so the array $(\pi_i(j))$ is of the kind counted by
$T(a,b)$, and the subarray with rows $I_{jj'}$ and columns $C_j$ is an
$a\times a$ Latin square on the symbol set $S_{j'}$.  Thus $L$ is
determined by the array together with one Latin square of order $a$ for
each of the $b^2$ pairs $(j,j')$, on the row set $I_{jj'}$, column set
$C_j$ and symbol set $S_{j'}$ determined by the array (identifying each
of these sets with $[a]$ in increasing order).  Conversely every such
array together with any $b^2$ small squares gives a square of the
required kind: row $i$ carries, on the column block $C_j$, each symbol
of $S_{\pi_i(j)}$ once, and these $b$ symbol blocks are distinct; column
$c\in C_j$ carries, on the rows $I_{jj'}$, each symbol of $S_{j'}$ once,
and the $b$ row sets $I_{jj'}$ ($j'\in[b]$) partition $[n]$.  There are
at most $(b!)^n$ arrays and $L(a)$ choices for each small square.
\end{proof}

\begin{proof}[Proof of Theorem~\ref{thm:main}]
Throughout, $\Pr$ denotes proportion in $\LL_n^0$; by the discussion above,
these proportions equal the corresponding proportions of reduced squares, hence of loops.
Let $L^0$ be uniform on $\LL_n^0$ and $H=H(L^0)$. Since $H$ is
transitive, if $H\not\supseteq A_n$ then $H$ is imprimitive or
primitive of order at most $B:=e^{4\sqrt n\log^2n}$ (for $n\ge n_0$, by
the group-theoretic facts above).

\emph{Imprimitive.}  Suppose $H$ preserves a partition $\Pi$ of $[n]$
into $b$ blocks of size $a$, $1<a<n$.  Then every row of $L^0$ maps
each block of $\Pi$ onto a block of $\Pi$, so by Lemma~\ref{lem:block}
with $\Pi'=\Pi$,
\[
\Pr\bigl[H\le\Stab(\Pi)\bigr]\le\frac{(b!)^n\,L(a)^{b^2}}{L(n)/n!}.
\]
There are $n!/((a!)^bb!)\le n!$ partitions $\Pi$ with blocks of size $a$
and fewer than $n$ divisors $a$ of $n$.  By \eqref{eq:L}, with $b=n/a$,
$(b!)^n\le b^{nb}$ and $\log n!\le n\log n$ (of the two factors $n!$
below, one is the $n!$ of $|\LL_n^0|=L(n)/n!$ and the other bounds the
number of partitions),
\begin{align*}
\log\frac{n!\cdot n!\,(b!)^n\,L(a)^{b^2}}{L(n)}
&\le2n\log n+\frac{n^2}a\log b+\Bigl(n^2\log a-2n^2+\frac{n^2}a\bigl(\tfrac12\log^2a+2\log a+3\bigr)\Bigr)\\
&\qquad-n^2\log n+2n^2\\
&=-n^2\log b+\frac{n^2}a\Bigl(\log b+\tfrac12\log^2a+2\log a+3\Bigr)+2n\log n.
\end{align*}
If $a\ge\sqrt n$, the middle term is at most
$n^{3/2}(\frac12\log^2n+3\log n+3)=o(n^2)$, so the right side is
$-n^2\log b+o(n^2)\le-n^2\log2+o(n^2)$.  If $2\le a<\sqrt n$, then
$b>\sqrt n$, $(n^2/a)\log b\le\frac12n^2\log b$, and
$(\frac12\log^2a+2\log a+3)/a\le\frac52$ (the function is decreasing in
$a\ge1$, its derivative being $-(\frac12\log^2a+\log a+1)/a^2$, and is
less than $2.32$ at $a=2$), so the right side is at most
$-\frac12n^2\log b+\frac52n^2+2n\log n\le-\frac14n^2\log n+3n^2$, which
is at most $-n^2\log2$ once $n\ge e^{15}$.  Summing over $\Pi$ and $a$,
\[
\Pr[H\text{ imprimitive}]\le n\,e^{-(\log2-o(1))n^2},
\]
where the $o(1)$ is $O(\log^2n/n)$: the term $a=n/2$ is at most
$e^{-n^2\log2+O(n\log^2n)}$, every term with $b\ge3$ and
$a\ge\sqrt n$ is at most $e^{-n^2\log3+O(n^{3/2}\log^2n)}$, and every
term with $a<\sqrt n$ is at most $e^{-n^2\log2}$ for $n\ge e^{15}$, so
the sum is at most $n\,e^{-n^2\log2+O(n\log^2n)}$.

\emph{Primitive.}  Let $n\ge n_0$.  For a fixed subgroup $K\le S_n$,
the squares $L^0\in\LL_n^0$ all of whose rows lie in $K$ number at most
$|K|^{n-1}$.  A subgroup of order at most $B$ is generated by at most
$\log_2B\le6\sqrt n\log^2n$ elements, so there are at most
$(n!)^{6\sqrt n\log^2n}\le e^{6n^{3/2}\log^3n}$ of them.  Summing over
these,
\begin{align*}
\Pr\bigl[H\text{ primitive},\ H\not\supseteq A_n\bigr]
&\le\frac{e^{6n^{3/2}\log^3n}\,B^{\,n-1}}{L(n)/n!}\\
&\le\exp\bigl(-n^2\log n+2n^2+n\log n+6n^{3/2}\log^3n+4n^{3/2}\log^2n\bigr)\\
&=e^{-(1-o(1))n^2\log n}.
\end{align*}
Adding the two cases gives the first inequality of the theorem.

\emph{The sign.}  $H\le A_n$ if and only if every row $\sigma_i$ of
$L^0$ is even.  Under the $n!$-to-one map $L\mapsto L^0$ this is the
condition that every row $r_i$ of $L\in\LL_n$ has the parity of $r_1$,
so the proportion of $L^0\in\LL_n^0$ with all rows even equals the
proportion of $L\in\LL_n$ with all rows of one parity.  Relabelling the
symbols by a transposition is a bijection of $\LL_n$ reversing the
parity of every row, so the squares with all rows odd are as many as
those with all rows even, and the latter are at most $c^n|\LL_n|$ by
H\"aggkvist and Janssen \cite{HJ96}.  Hence $\Pr[H\le A_n]\le2c^n$.  If $H\supseteq A_n$ and
$H\not\le A_n$ then $H=S_n$; so $\Pr[H\ne S_n]\le e^{-(\log2-o(1))n^2}+2c^n\le3c^n$
for large $n$.  Finally $\Mlt(Q)\supseteq\Mlt_\lambda(Q)=S_n$ is
$2$-transitive, which is the rank-$2$ condition of \cite{Smith}.
\end{proof}

\section{Remarks}

\begin{remark}[sharpness]\label{rem:sharp}
The first bound of the theorem is sharp on the scale $n^2$ for even
$n$ and not otherwise; more precisely, the proof and the lemma together
determine the exponent at every order.  Let $n$ be composite with least
prime factor $p$ (so $p\le\sqrt n$).  Then
\[
\log\Pr[H\not\supseteq A_n]=-n^2\log p+o(n^2),
\]
so that $\log\Pr[H\not\supseteq A_n]/n^2\to-\log2$ along even $n$ (no
statement about a multiplicative constant is intended) and the exponent
$\log2$ of the theorem is attained there, while for odd composite $n$
the true exponent is $\log p>\log2$.  For the lower bound take the
partition $\Pi$ into $p$ blocks of size $a=n/p$: $T(a,p)\ge1$ (the
array in which the rows of the $k$-th block of $a$ rows are the $k$-th
row of the cyclic Latin square of order $p$ has every symbol $a$ times
in every column), so by Lemma~\ref{lem:block} at least $L(a)^{p^2}$
squares of $\LL_n$ have all rows in $\Stab(\Pi)$, and by \eqref{eq:K}
and \eqref{eq:L}
\[
\Pr[H\not\supseteq A_n]\ge\Pr[H\le\Stab(\Pi)]\ge\frac{L(a)^{p^2}}{L(n)}
\ge e^{\,p^2(a^2\log a-2a^2)-n^2\log n+2n^2-n(\frac12\log^2n+2\log n+3)}
=e^{-n^2\log p-O(n\log^2n)}.
\]
For the upper bound, split the imprimitive sum at $a=\log^2n$ instead
of $\sqrt n$.  Every divisor $b=n/a>1$ is at least $p$.  If
$a\ge\log^2n$, the middle term of the display is at most
$n^2\log n/\log^2n+n^2(\frac12\log^2a+2\log a+3)/a=O(n^2/\log n)$
(the function of $a$ is decreasing), so the term is at most
$-n^2\log b+O(n^2/\log n)\le-n^2\log p+O(n^2/\log n)$.  If $2\le a<\log^2n$,
the bound $-\frac14n^2\log n+3n^2$ of the proof applies (for large
$n$, so that $\log^2n<\sqrt n$), and it is below $-n^2\log p$ once
$\frac14\log n-3\ge2\log\log n$, because $\log p\le\log a<2\log\log n$.  The
primitive case contributes $e^{-(1-o(1))n^2\log n}\le e^{-n^2\log p-(\frac12-o(1))n^2\log n}$.
Hence $\Pr[H\not\supseteq A_n]\le e^{-n^2\log p+O(n^2/\log n)}$ for
every composite $n$, without the restriction on $p$ that the case
split at $\sqrt n$ would impose.  For even $n$ the error terms are
sharper: $T(n/2,2)=\binom n{n/2}$ (an array with $b=2$ is determined by
the set of rows that swap the blocks, which has size $n/2$), so exactly
$\binom n{n/2}L(n/2)^4$ squares of $\LL_n$ have all rows in
$\Stab(\Pi)$, and by \eqref{eq:K} with $n!/|\Stab(\Pi)|=\frac12\binom n{n/2}$
and the asymptotics $\log L(m)=m^2\log m-2m^2+O(m\log^2m)$ that
\eqref{eq:L} encodes,
$\Pr[H\le\Stab(\Pi)]=2^{-n^2+O(n\log^2n)}$; the union over the
$\frac12\binom n{n/2}$ partitions into two blocks changes this by
$2^{O(n)}$, and the proof bounds the whole sum by $n\,e^{-n^2\log2+O(n\log^2n)}$,
so $\log\Pr[H\not\supseteq A_n]=-n^2\log2+O(n\log^2n)$ for even $n$.

For prime $n$ every transitive group is primitive, and the primitive
display of the proof gives
$\Pr[H\not\supseteq A_n]\le e^{-n^2\log n+2n^2+O(n^{3/2}\log^3n)}$.
This too is sharp to two terms: the Cayley tables of the
$(n-1)!/|\mathrm{Aut}(\mathbb Z_n)|=(n-2)!$ loops on $[n]$ with identity
$1$ that are isomorphic to $\mathbb Z_n$ are reduced squares with $H$ the
regular cyclic group of order $n$, so by \eqref{eq:R} and \eqref{eq:L}
\[
\Pr[H\not\supseteq A_n]\ge\frac{(n-2)!}{R(n)}=\frac{(n-2)!\,n!\,(n-1)!}{L(n)}
\ge e^{-n^2\log n+2n^2-O(n\log^2n)},
\]
and $\log\Pr[H\not\supseteq A_n]=-n^2\log n+2n^2+O(n^{3/2}\log^3n)$
along the primes.

The threshold $n\ge e^{15}$ in the proof is an artefact of the case
split and has not been optimised: the sum over the divisors $a$ of $n$
of $n!\cdot n!\,(b!)^nL(a)^{b^2}/L(n)$ (the first display of the
imprimitive case times the crude count $n!$ of partitions, with
\eqref{eq:L} inserted for $L(a)$ and $L(n)$ and $n!$, $(b!)^n$ exact)
is already less than $1$ for every composite $n\ge93$ (and is
$e^{-0.88\,n^2\log2}$ at $n=1000$), and with the exact van der
Waerden and Br\'egman expressions in place of \eqref{eq:L} and the
exact number $n!/((a!)^bb!)$ of partitions in place of $n!$ it is less
than $1$ for every composite $n\ge15$ ($e^{-0.37\,n^2\log2}$ at
$n=30$, $e^{-0.70\,n^2\log2}$ at $n=100$).  These statements are
certified as follows \cite{repo}: both sums were evaluated for every
composite $n\le25000$, and for $n\ge22107$ the two inequalities of the
proof --- a term with $a\ge\sqrt n$ is at most
$-n^2\log2+n^{3/2}(\frac12\log^2n+3\log n+3)+2n\log n$ and a term with
$a<\sqrt n$ at most $-\frac14n^2\log n+\frac52n^2+2n\log n$, both
decreasing in $n$ after division by $n^2$ --- make every term of the
first sum, hence of the second, which is termwise smaller, less than
$1/n$.
\end{remark}

\begin{remark}[the sign]\label{rem:sign}
The second bound of the theorem rests on the bound of \cite{HJ96} for
the proportion of all-even squares, which is $(\sqrt{3/4}+o(1))^n$, so
$c$ can be taken to be any number above $\sqrt{3/4}$.  Kwan, Petrova
and Sawhney \cite{KPS25} have since shown that the vector of row
parities of a uniformly random Latin square behaves like a vector of
independent fair coin flips, in particular that the all-even proportion
is at most $(\frac12+o(1))^n$; and Greenhill, Kopfov\'a, Kwan and
Petrova \cite{GKKP26} have now proved the matching lower
bound: the all-even proportion is $(\frac12+o(1))^n$.  The sign
paragraph of the proof turns this into
$\Pr[H\le A_n]=2\cdot(\frac12+o(1))^n=(\frac12+o(1))^n$, and since
$\Pr[H\le A_n,\,H\ne A_n]\le\Pr[H\not\supseteq A_n]=e^{-\Omega(n^2)}$,
\[
\Pr[H=A_n]=(\tfrac12+o(1))^n\qquad\text{and}\qquad\Pr[H\ne S_n]=(\tfrac12+o(1))^n,
\]
both two-sided; so the exceptional loops are, to the exponential
order, exactly those with $\Mlt_\lambda(Q)=A_n$, and in the theorem
every $c>\frac12$ works and no $c<\frac12$ does.
\end{remark}

\begin{remark}[quasigroups]
With $G=\langle r_1,\dots,r_n\rangle$ in place of $H$, the same three
steps (Lemma~\ref{lem:block} with $\Pi'=\Pi$, the count $|K|^n$, and
\cite{HJ96}) reprove Cameron's theorem \cite{Cam92} without
\cite{LP93}, with the same bounds: $e^{-(\log2-o(1))n^2}$ for
$G\not\supseteq A_n$ and $3c^n$ for $G\ne S_n$.  The converse direction is not
available: the loop case does not follow from the quasigroup case,
because the symbol relabelling that produces a reduced square replaces
the row group $\langle r_i\rangle$ by $\langle r_1^{-1}r_i\rangle$.
\end{remark}

\begin{remark}[the primitive case]\label{rem:primitive}
The primitive case needs only that every primitive $K\not\supseteq A_n$
has $\log|K|=o(n)$: then the generator count makes the number of such
$K$ at most $(n!)^{o(n)}=e^{o(n^2\log n)}$, negligible against
$L(n)/n!$.  Babai's bound $B$ gives this with room to spare, and so
would the bound $\exp(O(n^{1/3}\log^{7/3}n))$ of Sun and Wilmes
\cite{SW15}, also proved without the classification, which holds for
every uniprimitive group except $S_m$ and $A_m$ acting on $2$-subsets
of $[m]$ ($n=\binom m2$) and the primitive groups between $A_m\times A_m$
and $S_m\wr S_2$ in product action ($n=m^2$) \cite[Corollary~1.3]{SW15};
these exceptions are Cameron groups, of order $e^{O(\sqrt n\log n)}$,
and the first of them shows that Babai's bound is sharp up to the power
of the logarithm ($S_m$ on $2$-subsets has order $m!=e^{\Theta(\sqrt n\log n)}$).
Praeger and Saxl's elementary bound $|K|<4^n$ \cite{PS80} is too weak
for the generator count (it gives $(n!)^{2n}$ subgroups).  That is why
a first version of this note combined it with Pyber's bound
$2^{(\frac16\log_224+o(1))n^2}$ on the number of subgroups of $S_n$
\cite{Pyb93} (recently sharpened to $2^{n^2/16+O(n^{3/2})}$ by
Roney-Dougal and Tracey \cite{RT25}); we have not checked whether that
count can be obtained without the classification of finite simple
groups.  Eberhard pointed out that the elementary order bounds of
Babai and Pyber used above make the primitive case self-contained and
visibly free of it.
\end{remark}

\begin{remark}[what is and is not used]
Nothing is used about the distribution of any single row.  The event
$H\le K$ is the intersection of the $n-1$ events $\sigma_i\in K$, and the
identity row, which defeats the one-row argument, contributes nothing
to it.  The crude count $|K|^{n-1}$ against $(n!)^{n-1}e^{-n^2+o(n^2)}$
already settles each individual $K$ of index at least
$e^{(1+\delta)n}$; the two maximal subgroups of smaller index, $A_n$
and $S_{n/2}\wr S_2$ (with their conjugates), are where the Latin
constraints inside $K$ have to be used, through \cite{HJ96} and
through Lemma~\ref{lem:block} respectively.  In particular the theorem
says nothing about the conjecture of \cite{CGW08} on the cycle type of
the second row, which concerns one row given the identity row, and
which remains open; the consequence recorded in \cite[Lemma~6.2]{CGW08}
simply does not need it.
\end{remark}

\begin{remark}[Eberhard's sketch]
In comments on \cite{Cam15} (October 2016), Eberhard observed that
the number of Latin squares with all rows in $K$ is at most
$|K|^n=(n!)^n/[S_n:K]^n$ against $L(n)=(n!)^n/e^{n^2+o(n^2)}$, concluded
that only $A_n$ and $S_{n/2}\wr S_2$ need separate treatment, and
asked whether the method of \cite{HJ96} adapts to other subgroups;
Cameron replied that this ``should not be too hard'' and that the
second-row conjecture was the most interesting part; Eberhard then
sketched for $S_{n/2}\wr S_2$ a row-by-row count (fill the first
$k\approx n/10$ rows inside the group, bound the rest by the standard
row-by-row estimate).
The sketch is correct once one passes to maximal subgroups, which
takes care of the imprimitive groups with $b\ge3$; what remains to be
supplied is the primitive case, for which an elementary order bound for
primitive groups suffices, as Eberhard pointed out to us (the sketch's
``only $A_n$ and $S_{n/2}\wr S_2$'' presupposes such a bound; see
Remark~\ref{rem:primitive}), and the details for $S_{n/2}\wr S_2$.  Lemma~\ref{lem:block} treats all imprimitive groups
uniformly and gives the sharp exponent.  As far as we know the argument
was not written out.
\end{remark}

\begin{remark}[small orders]\label{rem:small}
For $n=5$, of the $56$ reduced squares $50$ have $H=S_5$ and $6$ have
$H=C_5$ (the sorted Cayley table of $\mathbb Z_5$ for each of the six
cyclic subgroups; these are also the all-even ones).  For $n=6$, of
$9408$ reduced squares $7776$ have $H=S_6$, $72$ have $A_6$ and $1560$
have an imprimitive $H$ (block sizes $2$ and $3$, orders $6$ to $48$).
In samples of $3000$ Latin squares at each order $7\le n\le16$
produced by the Jacobson--Matthews chain \cite{JM96} (the chain on
proper and improper squares, started from the cyclic square, run for
$200t$ moves and then recorded at those times $200t+kt$, $k\ge1$, at
which the state is proper, improper checkpoints being skipped with the
grid kept, $t=20n^2$; the stationary distribution restricted to proper
states is uniform, and recording at fixed times rather than at the first
proper state after a prescribed number of moves avoids a stopping-time
bias, so the protocol has the right target distribution; but we have
not verified mixing, so these are samples from a chain, not certified
uniform samples), the number of squares with $H=A_n$ was
$37,25,18,4,1,4,0,1,0,0$ for $n=7,\dots,16$, and every other $H$ was
$S_n$ except for two squares at $n=8$ with $H=S_2\wr S_4$ (order
$384$) \cite{repo}.
\end{remark}

\begin{remark}[labelled and unlabelled loops]
The theorem counts labelled loops on $[n]$ with identity $1$, i.e.\
reduced squares, and the error bound survives the passage to
isomorphism classes.  Let $p_n\le3c^n$ be the proportion of reduced
squares with $H\ne S_n$ and $\varepsilon_n$ the proportion of reduced
squares whose loop has a nontrivial automorphism group.  A loop
automorphism is an autotopism of the square, the autotopy group is an
isotopy invariant, and the map from $\LL_n$ to reduced squares is
uniform, so $\varepsilon_n$ is at most the proportion of Latin squares
with a nontrivial autoparatopy group, which is $n^{-3n^2/8+o(n^2)}$ by
\cite[Theorem~4]{MW05}.  An isomorphism class with trivial automorphism
group contains exactly $(n-1)!$ reduced squares, and every class
contains at least one, so the proportion of isomorphism classes with
multiplication group $\ne S_n$ is at most
$(p_n+(n-1)!\,\varepsilon_n)/(1-\varepsilon_n)=O(c^n)$: with $N=R(n)$,
the classes with $H\ne S_n$ and trivial automorphism group number at
most $p_nN/(n-1)!$ (the row group is an isomorphism invariant up to
conjugacy, so every square of such a class has $H\ne S_n$), the classes
with nontrivial automorphism group number at most $\varepsilon_nN$, and
the classes with trivial automorphism group alone number
$(1-\varepsilon_n)N/(n-1)!$.
\end{remark}

\paragraph{Acknowledgements.}  We thank Sean Eberhard for the argument
he sketched in 2016 and for pointing out that the elementary order
bounds of Babai and Pyber make the primitive case free of the
classification of finite simple groups.  [Cameron; Wanless.]

\begin{thebibliography}{9}
\bibitem{Bab81} L. Babai, On the order of uniprimitive permutation groups, \emph{Ann. of Math.} 113 (1981), 553--568.
\bibitem{Bab82} L. Babai, On the order of doubly transitive permutation groups, \emph{Invent. Math.} 65 (1982), 473--484.
\bibitem{Cam92} P.~J. Cameron, Almost all quasigroups have rank 2, \emph{Discrete Math.} 106/107 (1992), 111--115.
\bibitem{Cam15} P.~J. Cameron, A niggling problem, blog post, 24 January 2015, with comments by S.~Eberhard, 12--13 October 2016, \url{https://cameroncounts.wordpress.com/2015/01/24/a-niggling-problem/}.
\bibitem{CGW08} N.~J. Cavenagh, C. Greenhill and I.~M. Wanless, The cycle structure of two rows in a random Latin square, \emph{Random Structures Algorithms} 33 (2008), 286--309.
\bibitem{GKKP26} C. Greenhill, M. Kopfov\'a, M. Kwan and K. Petrova, On the probability that a random Latin square is all-even, arXiv:2610.13487 (2026).
\bibitem{HJ96} R. H\"aggkvist and J.~C.~M. Janssen, All-even Latin squares, \emph{Discrete Math.} 157 (1996), 199--206.
\bibitem{JM96} M.~T. Jacobson and P. Matthews, Generating uniformly distributed random Latin squares, \emph{J. Combin. Des.} 4 (1996), 405--437.
\bibitem{KPS25} M. Kwan, K. Petrova and M. Sawhney, Parities in random Latin squares, arXiv:2509.13125 (2025).
\bibitem{LP93} T. {\L}uczak and L. Pyber, On random generation of the symmetric group, \emph{Combin. Probab. Comput.} 2 (1993), 505--512.
\bibitem{MW05} B.~D. McKay and I.~M. Wanless, On the number of Latin squares, \emph{Ann. Comb.} 9 (2005), 335--344.
\bibitem{PS80} C.~E. Praeger and J. Saxl, On the orders of primitive permutation groups, \emph{Bull. London Math. Soc.} 12 (1980), 303--307.
\bibitem{Pyb93} L. Pyber, Enumerating finite groups of given order, \emph{Ann. of Math.} 137 (1993), 203--220.
\bibitem{Pyb93b} L. Pyber, On the orders of doubly transitive permutation groups, elementary estimates, \emph{J. Combin. Theory Ser.~A} 62 (1993), 361--366.
\bibitem{RT25} C.~M. Roney-Dougal and G. Tracey, Subgroups of symmetric groups: enumeration and asymptotic properties, arXiv:2503.05416 (2025).
\bibitem{repo} [author], \texttt{cgw-conjecture}, \url{https://github.com/mkkinyon/cgw-conjecture}, commit \texttt{1232610} (12 October 2026): the thresholds of Remark~\ref{rem:sharp} are the output of \texttt{python3 s13\_loopbounds.py 25000} (log \texttt{runs/s13/loops/bounds25000.log}); the samples of Remark~\ref{rem:small} are produced by \texttt{./jm \$n 3000 \$((20*n*n)) \$((100+n)) | python3 s13\_loopgroup.py sample \$n} (logs \texttt{runs/s13/loops/sample\$n.log}, each recording its command) and the exact counts at $n=5,6$ by \texttt{python3 s13\_loopgroup.py exact \$n} (\texttt{exact5.log}, \texttt{exact6.log}).
\bibitem{Smith} J.~D.~H. Smith, \emph{An Introduction to Quasigroups and Their Representations}, Chapman \& Hall/CRC, 2007.
\bibitem{SW15} X. Sun and J. Wilmes, Faster canonical forms for primitive coherent configurations, in \emph{STOC '15}, ACM, 2015, 693--702.
\bibitem{vLW} J.~H. van Lint and R.~M. Wilson, \emph{A Course in Combinatorics}, 2nd ed., Cambridge University Press, 2001.
\end{thebibliography}
\end{document}
